\begin{table}[!htbp]\centering\scriptsize\setlength{\tabcolsep}{3pt}
\caption{Les arbitrages pour des ménages types (variation de l'IR + CSG + CRDS, € par an)}\label{tab:ircsg_arbcas}
\resizebox{\linewidth}{!}{%
\begin{tabular}{@{}lrrrrr@{}}
\toprule
& Proposition & Quotient & Abatt. 10~\% & Chèque & Franchise 4\,000~€ \\
\midrule
Jeune salarié au Smic & $-$447 & $-$635 & $-$259 & $-$359 & $-$243 \\
Cadre seul, 100 000 € brut & $-$4\,015 & $-$2\,733 & $-$3\,682 & $-$4\,962 & $-$3\,894 \\
Parent isolé, deux enfants, au Smic & $-$2\,027 & $-$688 & $-$2\,027 & $-$2\,740 & $-$1\,683 \\
Couple sans enfant à un seul revenu, 80 000 € & +3\,498 & +65 & +3\,814 & +2\,873 & +3\,677 \\
Couple avec deux enfants à un seul revenu, 50 000 € & $-$862 & $-$77 & $-$539 & +92 & $-$518 \\
Couple avec trois enfants, 60 000 € brut (60/40) & $-$2\,922 & $-$993 & $-$2\,591 & $-$3\,078 & $-$2\,452 \\
Couple avec trois enfants à un seul revenu, 60 000 € & $-$1\,003 & +102 & $-$682 & +315 & $-$670 \\
Retraité seul, pension de 16 000 € & +1\,005 & +877 & +871 & +1\,041 & +840 \\
Retraité seul, pension de 20 000 € & +1\,353 & +1\,109 & +1\,262 & +1\,521 & +940 \\
Retraité seul, pension de 25 000 € & +1\,815 & +1\,577 & +1\,519 & +2\,118 & +1\,092 \\
Couple de retraités, 45 000 € de pensions (60/40) & +1\,330 & +196 & +809 & +1\,667 & +309 \\
\bottomrule
\end{tabular}}
\par\vspace{3pt}\parbox{\linewidth}{\raggedright\scriptsize\color{encre}
Mêmes cas-types que le tableau~\ref{tab:ircsg_cas}. Quotient : imposition commune des couples, quotient conjugal et familial aux plafonds actuels à la place de la franchise par enfant ; abattement de 10~\% sur les pensions en plus de l'exonération des petites pensions ; chèque de 1\,191~€ par enfant, de même coût que la franchise de 6\,000~€ par enfant ; franchise d'activité de 4\,000~€ par adulte au lieu de 8\,000~€. Chaque variante a son propre barème, recherché sous les mêmes contraintes.}
\end{table}
